% Generated draft: requires mathematical and source review.
\documentclass[11pt]{article}
\usepackage[margin=25mm]{geometry}
\usepackage{fontspec}
\setmainfont{Latin Modern Roman}
\usepackage{amsmath,amssymb,mathtools,unicode-math}
\setmathfont{Latin Modern Math}
\usepackage{enumitem,needspace,xurl,hyperref}
\setlength{\parindent}{0pt}
\setlength{\parskip}{5pt}
\setlength{\emergencystretch}{3em}
\newcommand{\N}{\mathbb{N}}
\newcommand{\Z}{\mathbb{Z}}
\newcommand{\Q}{\mathbb{Q}}
\newcommand{\R}{\mathbb{R}}
\newcommand{\C}{\mathbb{C}}
\newcommand{\F}{\mathbb{F}}
\newcommand{\A}{\mathbb{A}}
\newcommand{\PP}{\mathbb{P}}
\newcommand{\E}{\mathbb{E}}
\newcommand{\eps}{\varepsilon}
\newcommand{\dd}{\,\mathrm{d}}
\begin{document}
\section*{UnsolvedMath ID: 30005721}
\textbf{Source problem code:} catalogue entry
\subsection{Definitions and mathematical statement}
\textbf{Supplied question.} ``An important question is whether a pair of functions with low regularity can be multiplied.''

The supplied record does not define the domain of the functions, their scalar field, their regularity spaces or exponents, or the intended meaning of multiplication. In particular, it does not specify whether multiplication means pointwise multiplication, an extension of multiplication to distributions, or a renormalized operation; nor does it specify a target space or any required continuity, uniqueness, or approximation property. Consequently, no quantified mathematical assertion with definite hypotheses and conclusion can be recovered from this sentence alone.

\textbf{Unresolved source qualification.} The record labels the statement \texttt{unrecoverable}. Its supplied verification note says that the official report poses this question expositionally and immediately answers it using Young's theorem with an exact regularity criterion. That criterion and its hypotheses are absent from the supplied material, so neither the theorem nor a distinct remaining open problem can be stated here. The supplied \texttt{open} label must therefore be retained as a registry classification, not treated as a verified assertion that this extracted question is an open mathematical problem.
\subsection{Short English statement}
The source asks when functions of low regularity can be multiplied, but supplies no domains, regularity hypotheses, or definition of the product. Its verification note describes an expository question answered by Young's theorem. The missing criterion prevents recovery of a precise theorem or a definite open problem; this statement-only draft requires human mathematical review.
\subsection{Sources}
\begin{itemize}
\item Local UnsolvedMath registry, \texttt{data/unsolvedmath/problems.json}, permanent ID \(30005721\), title \emph{Multiplication of Low-Regularity Functions}, source code \texttt{catalogue entry}; supplied complete source record. Collection: \emph{Oberwolfach Reports Open Problems}. No individual proposer is identified; the supplied proposed year is 2023.
\item Supplied report citation: \emph{Arbeitsgemeinschaft: QFT and Stochastic PDEs (2024)}, Oberwolfach Reports, DOI 10.4171/owr/2023/59, pp.\ 3319--3370. Supplied source URL: \url{https://doi.org/10.4171/owr/2023/59}. Oberwolfach Reports are described in the record as documenting mathematical workshops and their open-problem discussions. The report citation is supplied provenance and has not been independently consulted here.
\item Supplied literature source: \emph{Oberwolfach Report 59/2023}, \url{https://doi.org/10.4171/owr/2023/59}. Literature assessment dated exactly \texttt{2026-08-22}: the supplied classification is \texttt{open}; the assessment says that the extracted sentence omits regularity spaces and the product notion, and that standard Young, paraproduct, and renormalization theories address restricted formulations rather than this unqualified question. This is supplied dated curation, described as literature triage derived from the supplied OWR curation, not an independently verified literature assessment. Its status and cited claims require independent verification.
\item Supplied statement-verification note: the official report asks the question expositionally and immediately answers it with Young's theorem, including an exact regularity criterion. The record does not reproduce that criterion. This supplied verification claim has not been independently checked here.
\end{itemize}
\end{document}
